\documentclass{article}


\usepackage{arxiv}

\usepackage[utf8]{inputenc} % allow utf-8 input
\usepackage[T1]{fontenc}    % use 8-bit T1 fonts
\usepackage{hyperref}       % hyperlinks
\usepackage{url}            % simple URL typesetting
\usepackage{booktabs}       % professional-quality tables
\usepackage{amsfonts}       % blackboard math symbols
\usepackage{nicefrac}       % compact symbols for 1/2, etc.
\usepackage{microtype}      % microtypography
\usepackage{lipsum}
\usepackage{graphicx}
\usepackage{amsmath} 
\usepackage{amssymb} % For symbols like \mathbb
\usepackage{amsthm} % For theorem-like environments
\usepackage{algorithm}
\usepackage{algpseudocode}

\graphicspath{ {./images/} }

\DeclareMathOperator*{\argmax}{argmax}

\title{Maximizing Algebraic Connectivity}


\author{
Author1 \\
  School of Coumputing and Information\\
  University of Pittsburgh\\
  Pittsburgh, PA 15213 \\
  \texttt{ziq2@pitt.edu} \\
  %% examples of more authors
   \And
Author2 \\
  School of Coumputing and Information\\
  University of Pittsburgh\\
  Pittsburgh, PA 15213 \\
  \texttt{ZIL50@pitt.edu} \\
  \And
Author3 \\
  School of Coumputing and Information\\
  University of Pittsburgh\\
  Pittsburgh, PA 15213 \\
  \texttt{yul217@pitt.edu} \\
  %% \AND
  %% Coauthor \\
  %% Affiliation \\
  %% Address \\
  %% \texttt{email} \\
  %% \And
  %% Coauthor \\
  %% Affiliation \\
  %% Address \\
  %% \texttt{email} \\
  %% \And
  %% Coauthor \\
  %% Affiliation \\
  %% Address \\
  %% \texttt{email} \\
}

\begin{document}
\maketitle
\begin{abstract}
Abstract
\end{abstract}


% keywords can be removed
%\keywords{First keyword \and Second keyword \and More}


\section{Introduction}

From power grids and communication backbones to the intricate web of social interactions, our world is defined by interconnected systems. Understanding the robustness of these networks, their ability to remain functional and cohesive, is a critical scientific and engineering challenge \cite{watts1998collective}. A powerful mathematical tool for this analysis begins with defining a graph. Let $G = (V, E)$ be a simple, undirected, and unweighted graph on $n = |V|$ vertices. Its structure can be represented by the graph Laplacian, $\mathbf{L} = \mathbf{D} - \mathbf{A}$, where $\mathbf{D}$ is the diagonal degree matrix and $\mathbf{A}$ is the adjacency matrix. The eigenvalues of this matrix, which are real and non-negative, encode deep structural properties of the graph \cite{biyikoglu2007laplacian}. The second-smallest eigenvalue, $\lambda_2$, famously named the \textit{algebraic connectivity} by Fiedler \cite{fiedler1973algebraic}, has proven to be of particular importance. It arises naturally in the analysis of dynamic processes on graphs, where its magnitude directly governs the convergence rate of consensus algorithms, the stability of synchronized systems, and the difficulty of partitioning the network.

Given this direct link between $\lambda_2$ and network performance, a central problem in network design is the $(n, m)$-maximization problem: construct a graph $G^*$ from the set of all simple, undirected graphs $\mathcal{G}(n, m)$ with $n$ nodes and $m$ edges, such that $\lambda_2(G^*) \ge \lambda_2(G)$ for all $G \in \mathcal{G}(n, m)$. This is, in effect, the search for the most robust graph given a fixed budget of nodes and edges. However, the problem is notoriously difficult. The search space $\mathcal{G}(n, m)$ is combinatorially vast, and the objective function $\lambda_2$ is a complex, non-convex function of the discrete edge set. The problem is known to be NP-hard \cite{das2004max}, which prohibits brute-force or exhaustive search for any practical network size and demands more sophisticated optimization strategies.

Over the past few decades, a wide array of methods has been proposed to tackle this challenge, generally falling into two distinct camps. On one side, powerful optimization techniques such as Semidefinite Programming (SDP) relaxation can find high-quality solutions by convexifying the non-convex search space \cite{ghosh2006growing, de2013semidefinite}. However, SDP-based methods suffer from significant scalability issues, with runtimes that make them impractical for the large-scale networks often encountered in real-world applications. On the other end of the spectrum, fast heuristics are common. These often involve greedily adding edges based on spectral or resistance-based metrics. Prominent examples include adding edges to maximize the quadratic form $x^T \mathbf{L} x$ using the Fiedler vector \cite{ghosh2006growing} or minimizing metrics like the total effective resistance \cite{ghosh2008maximizing}. While scalable, these greedy methods frequently get trapped in poor local optima. Recent work continues to explore the complex optimization landscape, identifying the properties of local and global maximizer graphs \cite{belk2023algebraic}.

This trade-off, between computationally prohibitive, high-quality solutions and fast, suboptimal heuristics motivates our work. In this paper, we propose a novel two-phase hybrid framework that bridges this gap. Our approach leverages the strengths of both deterministic construction and intelligent search. In Phase 1, we develop a novel $O(N^2)$ algorithm to construct a high-connectivity $k$-regular "skeleton" graph. This deterministic phase builds a strong structural foundation, ensuring the graph is well-connected from the start. In Phase 2, we employ a Reinforcement Learning (RL) agent to intelligently add the remaining $m - (nk/2)$ residual edges. We formulate this as a sequential decision problem, where the agent learns a policy to select edge additions that maximize the change in $\lambda_2$. This use of RL is inspired by its recent success in solving other hard combinatorial optimization problems \cite{bello2016neural}.

We hypothesize that this hybrid method, using a fast algorithm to get close and a smart RL agent to perform the final, difficult optimization, will produce graphs with higher algebraic connectivity than existing methods. We demonstrate through intensive computational experiments that our two-phase approach consistently outperforms a suite of well-known baselines, including SDP relaxation, Fiedler vector heuristics, and small-world models.

The remainder of this paper is organized as follows. Section 2 formally defines the $(n, m)$-maximization problem and reviews related work. Section 3 provides a high-level overview of our two-phase framework. Section 4 presents the novel $O(N^2)$ skeleton-building algorithm, and Section 5 details the Reinforcement Learning formulation for edge augmentation. Section 6 presents our comprehensive experimental validation and analysis. Finally, Section 7 concludes the paper and discusses avenues for future work.

\section{Problem Formulation and Related Work}

\subsection{The $(n, m)$-Maximization Problem}

Let $G = (V, E)$ be a simple, undirected graph with $n = |V|$ vertices and $m = |E|$ edges. The Laplacian matrix $\mathbf{L} \in \mathbb{R}^{n \times n}$ is defined as $\mathbf{L} = \mathbf{D} - \mathbf{A}$, where $\mathbf{D} = \text{diag}(d_1, \dots, d_n)$ is the diagonal matrix of vertex degrees $d_i$, and $\mathbf{A}$ is the adjacency matrix. $\mathbf{L}$ is a positive semidefinite (PSD) matrix with eigenvalues $0 = \lambda_1 \le \lambda_2 \le \dots \le \lambda_n$. The first eigenvalue $\lambda_1$ is always zero, corresponding to the eigenvector $\mathbf{1}$ (the vector of all ones).

The algebraic connectivity, $\lambda_2$, is given by the Courant-Fischer theorem as:
\begin{equation}
    \lambda_2(G) = \min_{x \perp \mathbf{1}, x \neq 0} \frac{x^T \mathbf{L} x}{x^T x} = \min_{x \perp \mathbf{1}, \|x\|_2=1} x^T \mathbf{L} x
\end{equation}
The corresponding eigenvector, often called the Fiedler vector, provides a spectral approximation for partitioning the graph \cite{fiedler1973algebraic}. A key property for optimization is the monotonicity of $\lambda_2$. Adding an edge $e = (u, v)$ to $G$ to create $G' = (V, E \cup \{e\})$ results in a new Laplacian $\mathbf{L}' = \mathbf{L} + \mathbf{b}\mathbf{b}^T$, where $\mathbf{b}$ is a vector with $1$ at index $u$, $-1$ at index $v$, and $0$ elsewhere. This rank-one update is positive semidefinite, and by Weyl's inequality, all eigenvalues of $\mathbf{L}'$ are greater than or equal to those of $\mathbf{L}$ \cite{HornJohnson2013}. Thus, $\lambda_2(G') \ge \lambda_2(G)$, meaning that adding edges never decreases the algebraic connectivity.

This work addresses the $(n, m)$-maximization problem, which is formally stated as finding the optimal graph $G^*$ in the set $\mathcal{G}(n, m)$ of all simple graphs on $n$ vertices and $m$ edges:
\begin{equation}
    G^* = \argmax_{G \in \mathcal{G}(n, m)} \lambda_2(G)
\end{equation}
This problem is a specific instance of the more general graph augmentation problem, which seeks to add $k$ edges to an existing graph $G_0$ to maximize $\lambda_2(G_0 + k \text{ edges})$. This augmentation problem has been proven to be NP-hard \cite{MOSKAOYAMA2008677}, implying the difficulty of our primary problem.

\subsection{Review of Existing Maximization Strategies}

The difficulty of the $(n, m)$-maximization problem has led to a diverse range of strategies, which we categorize into four main groups.

\textbf{1. Heuristic and Greedy Algorithms.} These methods are prized for their scalability and are the most common in practice. A well-known method, the Fiedler vector heuristic \cite{ghosh2006growing}, involves greedily adding edges $(u, v)$ that maximize the potential increase in $\lambda_2$, often approximated by $(x_u - x_v)^2$, where $x$ is the Fiedler vector. A related approach greedily adds edges that minimize the total \textit{effective resistance} of the graph, which is also connected to the Laplacian \cite{doi:10.1137/050645452}. Other heuristics leverage different structural properties, such as adding edges based on vertex degree and graph distance \cite{8412478}. While fast, these methods are myopic and often converge to solutions that are only locally optimal.

\textbf{2. Optimization-Based Methods.} To find higher-quality solutions, many researchers have formulated the problem as a continuous optimization. The most prominent of these is the Semidefinite Programming (SDP) relaxation \cite{ghosh2006growing, de2013semidefinite}. By relaxing the discrete edge variables to be continuous and imposing a set of PSD constraints, these methods can often find near-global optima. However, the computational cost of solving SDPs, typically scaling poorly with $n$, renders them unsuitable for large-scale graphs, though they serve as a crucial benchmark for solution quality.

\textbf{3. Structural and Algorithmic Approaches.} A third line of inquiry investigates the properties of optimal graphs to guide their construction. This includes theoretical work on the characteristics of local and global maximizers \cite{belk2023algebraic, 10004730}. For the specific case of $k$-regular graphs (where $m = nk/2$), algorithms have been developed to explicitly search for maximizers. A recent example includes a Depth-First Search (DFS) algorithm with pruning to find optimal $k$-regular graphs \cite{10818309}. Our Phase 1 skeleton-building algorithm is motivated by this line of work, but is designed for $O(n^2)$ performance rather than exhaustive search.

\textbf{4. Proxy-Based Optimization.} A distinct approach avoids the direct optimization of $\lambda_2$ altogether. Instead, it focuses on optimizing a related, but computationally simpler, spectral metric that also correlates with fast consensus. A prime example is the minimization of \textit{Laplacian Energy} (the sum of square of Laplacian eigenvalues), which can be solved more efficiently and has been shown to produce graphs that also possess large algebraic connectivity \cite{lu2024fast}.

\textbf{Baselines and Machine Learning.} Our work is benchmarked against methods from the first two categories, as well as against standard graph models like the small-world network \cite{watts1998collective}, which is known to have good connectivity and fast consensus properties but is not explicitly optimized for $\lambda_2$ \cite{1470321}. Finally, the use of Reinforcement Learning for graph problems is a nascent but promising field. RL has been successfully applied to general combinatorial optimization \cite{bello2016neural} and, more specifically, to goal-directed graph construction \cite{doi:10.1098/rspa.2021.0168}. Our work builds on this premise, using RL not to build a graph from scratch, but to learn an intelligent \emph{rewiring} policy that refines a deterministic initialization.




\section{REFINE: Rewiring via Edge-level Fiedler INferencE}

The core idea behind REFINE is to cast the $(n,m)$-maximization problem as a \emph{graph rewiring} Markov Decision Process (MDP), where a learned policy iteratively adds and removes edges to improve $\lambda_2$ while preserving the edge budget $m$. Rather than constructing a graph from scratch—a task with an overwhelmingly large action space of $O\!\left(\binom{n(n-1)/2}{m}\right)$—we initialize with a simple connected graph and refine it through local edge swaps. This section presents the state representation, the policy architecture, the subspace tracking mechanism that enables efficient inference, and the training procedure.

\subsection{MDP Formulation}

We define the rewiring MDP $\mathcal{M} = (\mathcal{S}, \mathcal{A}, \mathcal{T}, \mathcal{R}, \gamma)$ as follows.

\textbf{State.} A state $s \in \mathcal{S}$ is a graph $G = (V, E)$ with $|V| = n$ and $|E| = m$, together with a tracked spectral subspace (defined in Section~\ref{sec:rr}). The initial state $s_0$ is constructed by forming a Hamiltonian cycle on $n$ nodes and adding $m - n$ edges uniformly at random, ensuring connectivity.

\textbf{Action.} Each action $a = (e^+, e^-)$ is a \emph{swap}: adding a non-edge $e^+ \in \binom{V}{2} \setminus E$ and removing an existing non-bridge edge $e^- \in E$. Bridge edges—those whose removal disconnects the graph—are excluded from the removal set, which is computed via Tarjan's algorithm in $O(n + m)$ time. This ensures $|E| = m$ after every action and guarantees the graph remains connected.

\textbf{Transition.} An episode proceeds over $K$ \emph{rounds}. Within each round, the agent performs $\Delta = \lfloor \alpha m \rfloor$ sequential swaps, where $\alpha$ is a swap fraction (typically $\alpha = 0.05$). Each round is further decomposed into two phases: $\Delta$ edge additions followed by $\Delta$ matched edge removals, yielding a net-zero change in edge count.

\textbf{Reward.} Let $\lambda_2^{(k)}$ denote the algebraic connectivity at the beginning of round $k$. After the ADD phase of round $k$, an exact eigensolve yields $\lambda_2^{\text{post-add}}$, and the ADD reward is
\begin{equation}\label{eq:reward_add}
    r_k^{\text{add}} = \lambda_2^{\text{post-add}} - \lambda_2^{(k)}.
\end{equation}
After the subsequent REM phase, a second eigensolve yields $\lambda_2^{(k+1)}$, and the REM reward is
\begin{equation}\label{eq:reward_rem}
    r_k^{\text{rem}} = \lambda_2^{(k+1)} - \lambda_2^{\text{post-add}}.
\end{equation}
Within each phase, only the \emph{last} action receives the nonzero reward; all preceding actions receive zero. Credit is then propagated backward via Generalized Advantage Estimation (GAE) \cite{schulman2015high}. This batch-reward scheme requires only two eigensolves per round (each $O(n^3)$) and avoids the per-swap myopia that arises from rewarding each action independently.

\subsection{Edge-Level Feature Representation}

A key design choice is that the policy operates on \emph{edge-level features} rather than node embeddings or adjacency-based representations. For every pair $(i, j) \in V \times V$, we define a feature vector $\phi_{ij} \in \mathbb{R}^4$:
\begin{equation}\label{eq:features}
    \phi_{ij} = \Bigl[\, n \cdot (v_{2,i} - v_{2,j})^2, \;\; \tfrac{d_i}{n-1}, \;\; \tfrac{d_j}{n-1}, \;\; \tfrac{k}{K} \,\Bigr]
\end{equation}
where $v_2$ is the (tracked) Fiedler vector, $d_i$ is the degree of node $i$, $k$ is the current round, and $K$ is the total number of rounds. The first component, the \emph{Fiedler gap}, measures the squared difference in the Fiedler vector components between nodes $i$ and $j$. This feature is directly motivated by the Courant-Fischer characterization: adding an edge $(u,v)$ increases $\lambda_2$ by approximately $(v_{2,u} - v_{2,v})^2$ in the first-order perturbation. An ablation study (Section~\ref{sec:ablation}) confirmed that this single spectral feature, combined with degree information, captures the essential signal for edge selection, and that additional eigenvector components ($v_3, v_4, v_5$) and spectral gap scalars serve primarily as noise dimensions that degrade out-of-distribution generalization.

For the value function, we define graph-level features $\psi \in \mathbb{R}^3$:
\begin{equation}\label{eq:graph_features}
    \psi = \Bigl[\, \tfrac{\bar{d}}{n-1}, \;\; \tfrac{\lambda_2}{n}, \;\; \tfrac{k}{K} \,\Bigr]
\end{equation}
where $\bar{d}$ is the mean degree. All features are normalized by graph size, enabling a policy trained on small graphs ($n \leq 16$) to generalize to much larger instances.

\subsection{Policy Architecture}

The policy consists of three independent multi-layer perceptrons (MLPs), totaling approximately $13{,}600$ parameters:

\begin{itemize}
    \item \textbf{ADD network} $\pi_{\text{add}}$: maps $\phi_{ij} \in \mathbb{R}^4$ to a scalar logit $\ell_{ij}^+$ for adding non-edge $(i,j)$,
    \item \textbf{REM network} $\pi_{\text{rem}}$: maps $\phi_{ij} \in \mathbb{R}^4$ to a scalar logit $\ell_{ij}^-$ for removing edge $(i,j)$,
    \item \textbf{Value network} $V_\theta$: maps $\psi \in \mathbb{R}^3$ to a scalar state-value estimate $\hat{V}(s)$.
\end{itemize}
Each MLP has the architecture $\text{Linear}(d_{\text{in}}, 64) \to \text{SiLU} \to \text{Linear}(64, 64) \to \text{SiLU} \to \text{Linear}(64, 1)$, with orthogonal weight initialization (gain $= 0.5$). The ADD and REM networks have independent weights, allowing them to learn distinct criteria for edge addition versus removal.

During training, actions are sampled from the softmax distribution over valid actions:
\begin{equation}\label{eq:policy}
    \pi(e^+ \mid s) = \frac{\exp(\ell_{ij}^+)}{\sum_{(i',j') \in \mathcal{A}^+} \exp(\ell_{i'j'}^+)}, \quad
    \pi(e^- \mid s) = \frac{\exp(\ell_{ij}^-)}{\sum_{(i',j') \in \mathcal{A}^-} \exp(\ell_{i'j'}^-)}
\end{equation}
where $\mathcal{A}^+$ is the set of non-edges (upper triangular) and $\mathcal{A}^-$ is the set of non-bridge edges excluding those added in the current round. At inference, actions are selected greedily via $\arg\max$.

\subsection{Rayleigh-Ritz Subspace Tracking}\label{sec:rr}

A central challenge is that the Fiedler vector $v_2$—the dominant feature—changes after every edge swap. Recomputing it via full eigendecomposition costs $O(n^3)$ per swap, which is prohibitive at inference for large $n$. We address this with a Rayleigh-Ritz (RR) subspace tracking scheme that maintains an approximate spectral decomposition in $O(nk^2)$ time per swap, where $k \ll n$ is the subspace dimension.

We track a $k$-dimensional subspace $\mathbf{V} = [v_2, \ldots, v_{k+1}] \in \mathbb{R}^{n \times k}$ spanning the bottom $k$ nontrivial eigenvectors of the Laplacian, together with the corresponding eigenvalues $\boldsymbol{\lambda} = [\lambda_2, \ldots, \lambda_{k+1}]$. When an edge $(u, v)$ is added or removed, the Laplacian undergoes a rank-one update:
\begin{equation}\label{eq:laplacian_update}
    \mathbf{L}' = \mathbf{L} \pm \mathbf{z}\mathbf{z}^T, \quad \mathbf{z} = \mathbf{e}_u - \mathbf{e}_v
\end{equation}
where $\mathbf{e}_i$ is the $i$-th standard basis vector. The RR update projects this perturbation into the tracked subspace:

\begin{enumerate}
    \item Compute the gap vector $\boldsymbol{\delta} \in \mathbb{R}^k$, where $\delta_\ell = V_{u\ell} - V_{v\ell}$ for $\ell = 1, \ldots, k$. This costs $O(k)$.
    \item Form the projected Laplacian $\mathbf{L}_{\text{sub}} = \text{diag}(\boldsymbol{\lambda}) \pm \boldsymbol{\delta}\boldsymbol{\delta}^T \in \mathbb{R}^{k \times k}$. This costs $O(k^2)$.
    \item Diagonalize: $\mathbf{L}_{\text{sub}} = \mathbf{R} \boldsymbol{\Lambda}' \mathbf{R}^T$, yielding updated eigenvalues $\boldsymbol{\Lambda}'$ and a rotation matrix $\mathbf{R} \in \mathbb{R}^{k \times k}$. This costs $O(k^3)$.
    \item Rotate the full subspace: $\mathbf{V}' = \mathbf{V} \mathbf{R}$, costing $O(nk^2)$.
\end{enumerate}
For $k = 8$ (used throughout), each RR update costs $O(64n)$—effectively linear in $n$. At the beginning of each of the $K$ rounds, the subspace is re-initialized from a warm-started Lanczos iteration in $O(n^2)$ time, which corrects any accumulated drift from the rank-one projections.

\subsection{Inference Algorithm}

Algorithm~\ref{alg:refine_inference} presents the complete inference procedure. The key property is that \emph{no eigensolves are performed}: the Fiedler vector is maintained entirely through Lanczos initialization and RR tracking. After the initial batch scoring in $O(n^2)$, each swap requires only $O(n)$ work to re-score the affected rows and columns, yielding an overall complexity of $O(n^3)$ for dense graphs—matching the baseline heuristics.

\begin{algorithm}[t]
\caption{REFINE Inference}\label{alg:refine_inference}
\begin{algorithmic}[1]
\Require Graph size $n$, edge budget $m$, trained policy $(\pi_{\text{add}}, \pi_{\text{rem}})$, rounds $K$, swap fraction $\alpha$
\Ensure Graph $G = (V, E)$ with $|E| = m$ and improved $\lambda_2$
\State Initialize $G \gets \text{Ring}(n) \cup \text{RandomEdges}(m - n)$
\State $\Delta \gets \lfloor \alpha \cdot m \rfloor$
\For{$k = 0, \ldots, K-1$}
    \State $\mathbf{V}, \boldsymbol{\lambda} \gets \textsc{Lanczos}(\mathbf{L}_G, k{=}8)$ \Comment{$O(n^2)$, warm-started}
    \Statex \hspace{\algorithmicindent}\textbf{--- ADD Phase ---}
    \State Compute $\phi_{ij}$ for all $(i,j)$; \, $\boldsymbol{\ell}^+ \gets \pi_{\text{add}}(\boldsymbol{\phi})$ \Comment{$O(n^2)$}
    \For{$s = 1, \ldots, \Delta$}
        \State $(u, v) \gets \arg\max_{(i,j) \notin E} \ell_{ij}^+$
        \State $E \gets E \cup \{(u,v)\}$; \, update degrees
        \State $\mathbf{V}, \boldsymbol{\lambda} \gets \textsc{RR-Update}(\mathbf{V}, \boldsymbol{\lambda}, u, v, {+}1)$ \Comment{$O(nk^2)$}
        \State Re-score rows/columns $u, v$ in $\boldsymbol{\ell}^+$ \Comment{$O(n)$}
    \EndFor
    \Statex \hspace{\algorithmicindent}\textbf{--- REM Phase ---}
    \State $\mathbf{V}, \boldsymbol{\lambda} \gets \textsc{Lanczos}(\mathbf{L}_G, k{=}8)$ \Comment{Re-init RR}
    \State Compute $\phi_{ij}$ for all $(i,j)$; \, $\boldsymbol{\ell}^- \gets \pi_{\text{rem}}(\boldsymbol{\phi})$ \Comment{$O(n^2)$}
    \For{$s = 1, \ldots, \Delta$}
        \State $\mathcal{B} \gets \textsc{Bridges}(G)$ \Comment{$O(n + m)$}
        \State $(u, v) \gets \arg\max_{(i,j) \in E \setminus \mathcal{B}} \ell_{ij}^-$
        \State $E \gets E \setminus \{(u,v)\}$; \, update degrees
        \State $\mathbf{V}, \boldsymbol{\lambda} \gets \textsc{RR-Update}(\mathbf{V}, \boldsymbol{\lambda}, u, v, {-}1)$ \Comment{$O(nk^2)$}
        \State Re-score rows/columns $u, v$ in $\boldsymbol{\ell}^-$ \Comment{$O(n)$}
    \EndFor
\EndFor
\State \Return $G$
\end{algorithmic}
\end{algorithm}

\subsection{Training with Proximal Policy Optimization}

The policy is trained via Proximal Policy Optimization (PPO) \cite{schulman2017proximal} on graph instances sampled uniformly with $n \in [n_{\min}, n_{\max}]$ and $m \in [n-1, \binom{n}{2}]$. A critical design decision is the alignment between the data collection policy and the PPO re-evaluation policy: at each swap within a round, the agent observes the \emph{current} spectral state (rebuilt features from the current RR subspace), and these per-swap feature snapshots are stored alongside the corresponding action and log-probability. During PPO updates, the policy is re-evaluated on the \emph{exact} feature tensor the agent observed when it selected each action, ensuring correct policy gradients.

Algorithm~\ref{alg:refine_train} summarizes the training loop. For a batch of $B$ episodes, we collect transitions, compute GAE advantages with discount $\gamma$ and trace-decay $\lambda_{\text{GAE}}$, and perform $E$ epochs of clipped PPO updates. The entropy coefficient is annealed linearly from $\beta_{\text{start}}$ to $\beta_{\text{end}}$ over the course of training to encourage early exploration.

\begin{algorithm}[t]
\caption{REFINE Training}\label{alg:refine_train}
\begin{algorithmic}[1]
\Require Training range $[n_{\min}, n_{\max}]$, episodes $T$, batch size $B$, rounds $K$, swap fraction $\alpha$, PPO epochs $E$
\Ensure Trained policy parameters $\theta$
\State Initialize $\theta$ with orthogonal weights
\For{each PPO update $t = 1, \ldots, T/B$}
    \State $\mathcal{D}_{\text{add}}, \mathcal{D}_{\text{rem}} \gets \emptyset$
    \For{$b = 1, \ldots, B$} \Comment{Collect batch}
        \State Sample $n \sim \text{Uniform}[n_{\min}, n_{\max}]$, $m \sim \text{Uniform}[n{-}1, \binom{n}{2}]$
        \State Initialize $G$; \, $\Delta \gets \lfloor \alpha m \rfloor$
        \For{$k = 0, \ldots, K-1$}
            \State $\mathbf{V}, \boldsymbol{\lambda} \gets \textsc{Lanczos}(\mathbf{L}_G, k{=}8)$
            \State $\lambda_2^{\text{pre}} \gets \textsc{Eigensolve}(\mathbf{L}_G)$ \Comment{$O(n^3)$}
            \Statex \hspace{2\algorithmicindent}\textbf{--- ADD Phase ---}
            \For{$s = 1, \ldots, \Delta$}
                \State $\boldsymbol{\phi} \gets \textsc{BuildFeatures}(\mathbf{V}, \boldsymbol{\lambda}, G)$ \Comment{$O(n^2)$}
                \State $e^+ \sim \pi_{\text{add}}(\cdot \mid \boldsymbol{\phi})$ \Comment{Softmax sampling}
                \State Execute add; RR-update
                \State Store $(\boldsymbol{\phi}, e^+, \log\pi, \hat{V})$ in $\mathcal{D}_{\text{add}}$
            \EndFor
            \State $\lambda_2^{\text{post-add}} \gets \textsc{Eigensolve}(\mathbf{L}_G)$
            \State Assign $r^{\text{add}} = \lambda_2^{\text{post-add}} - \lambda_2^{\text{pre}}$ to last ADD transition
            \Statex \hspace{2\algorithmicindent}\textbf{--- REM Phase ---}
            \State $\mathbf{V}, \boldsymbol{\lambda} \gets \textsc{Lanczos}(\mathbf{L}_G, k{=}8)$ \Comment{Fresh RR}
            \For{$s = 1, \ldots, \Delta$}
                \State $\boldsymbol{\phi} \gets \textsc{BuildFeatures}(\mathbf{V}, \boldsymbol{\lambda}, G)$
                \State $e^- \sim \pi_{\text{rem}}(\cdot \mid \boldsymbol{\phi})$, $e^- \notin \mathcal{B}$
                \State Execute remove; RR-update
                \State Store $(\boldsymbol{\phi}, e^-, \log\pi, \hat{V})$ in $\mathcal{D}_{\text{rem}}$
            \EndFor
            \State $\lambda_2^{(k+1)} \gets \textsc{Eigensolve}(\mathbf{L}_G)$
            \State Assign $r^{\text{rem}} = \lambda_2^{(k+1)} - \lambda_2^{\text{post-add}}$ to last REM transition
        \EndFor
    \EndFor
    \State Compute GAE advantages for $\mathcal{D}_{\text{add}}$ and $\mathcal{D}_{\text{rem}}$ separately
    \For{$e = 1, \ldots, E$} \Comment{PPO update}
        \State Re-evaluate $\log\pi_\theta(e \mid \boldsymbol{\phi})$ and $V_\theta(\psi)$ on stored features
        \State $r(\theta) \gets \exp\bigl(\log\pi_\theta - \log\pi_{\text{old}}\bigr)$
        \State $\mathcal{L}_{\text{clip}} \gets -\mathbb{E}\bigl[\min\bigl(r(\theta)\hat{A},\; \text{clip}(r(\theta), 1{\pm}\epsilon)\hat{A}\bigr)\bigr]$
        \State $\mathcal{L}_{\text{val}} \gets \mathbb{E}\bigl[(V_\theta - R)^2\bigr]$
        \State $\theta \gets \theta - \eta \nabla_\theta \bigl(\mathcal{L}_{\text{clip}} + \tfrac{1}{2}\mathcal{L}_{\text{val}} - \beta \mathcal{H}[\pi_\theta]\bigr)$
    \EndFor
\EndFor
\end{algorithmic}
\end{algorithm}

\subsection{Complexity Analysis}

Table~\ref{tab:complexity} summarizes the per-episode complexity of REFINE. Let $\Delta = \lfloor \alpha m \rfloor$ denote the number of swaps per round. For dense graphs where $m = \Theta(n^2)$, we have $\Delta = \Theta(n^2)$, and the total number of swaps across $K$ rounds is $K\Delta = \Theta(n^2)$ (with $K$ constant).

\textbf{Inference.} At each round, one batch MLP evaluation scores all $O(n^2)$ pairs. Each of the $\Delta$ subsequent swaps performs a row/column re-score in $O(n)$ and an RR subspace update in $O(nk^2)$, both called $2K\Delta$ times across the episode. Bridge detection runs once per round in $O(n + m)$. Summing across all $K$ rounds: the row/column re-scoring contributes $O(nK\Delta) = O(n^3)$ for dense graphs, and the RR updates contribute $O(nk^2 K\Delta) = O(n^3)$ (since $k = 8$ is constant). No eigensolves are performed. The total inference complexity is therefore $O(n^3)$, matching the FV and ER baselines.

\textbf{Training} ($n \leq 16$). The key difference is that features are \emph{fully rebuilt} at every swap for PPO alignment, costing $O(n^2)$ per swap rather than $O(n)$ for row re-scoring. This yields $O(n^2 K\Delta) = O(n^4)$ for dense graphs. Additionally, two eigensolves ($O(n^3)$ each) are performed per round, contributing $O(n^3 K)$. At the training sizes $n \leq 16$, these costs are negligible: $16^4 \approx 65{,}000$ and $16^3 \cdot 20 \approx 82{,}000$.

\begin{table}[t]
\centering
\caption{Per-episode complexity of REFINE. $\Delta = \lfloor\alpha m\rfloor$ swaps per round, $K$ rounds per episode. For dense graphs, $\Delta = \Theta(n^2)$.}\label{tab:complexity}
\begin{tabular}{@{}lccc@{}}
\toprule
\textbf{Component} & \textbf{Per-call} & \textbf{Training total} & \textbf{Inference total} \\
\midrule
Lanczos init (warm-started) & $O(n^2)$ & $O(Kn^2)$ & $O(Kn^2)$ \\
Feature build / batch MLP & $O(n^2)$ & $O(n^2 K\Delta)$ & $O(Kn^2)$ \\
Row/col re-scoring & $O(n)$ & --- & $O(nK\Delta)$ \\
RR subspace update ($k{=}8$) & $O(nk^2)$ & $O(nk^2 K\Delta)$ & $O(nk^2 K\Delta)$ \\
Bridge detection (per round) & $O(n + m)$ & $O(K(n{+}m))$ & $O(K(n{+}m))$ \\
Eigensolve (per round) & $O(n^3)$ & $O(Kn^3)$ & \textbf{---} \\
\midrule
\textbf{Total (dense, $m = \Theta(n^2)$)} & & $\boldsymbol{O(n^4)}$ & $\boldsymbol{O(n^3)}$ \\
\bottomrule
\end{tabular}
\end{table}

\section{Comparative Analysis of Algebraic Connectivity}

\subsection{Baselines}\label{sec:baselines}

We compare REFINE against four established baselines, all evaluated on the same $(n, m)$ instances.

\textbf{Fiedler Vector Heuristic (FV).} Starting from a spanning tree, edges are greedily added to maximize the quadratic form $(v_{2,u} - v_{2,v})^2$, recomputing the Fiedler vector after each addition \cite{ghosh2006growing}. This is a deterministic $O(mn^3)$ procedure.

\textbf{Effective Resistance (ER).} Edges are greedily added to maximize the decrease in total effective resistance, which is equivalent to maximizing $R_{\text{eff}}(u, v)$ for the added edge \cite{ghosh2008maximizing}. This also requires $O(mn^3)$ time due to the per-edge eigendecomposition.

\textbf{Small-World (SW).} The Watts-Strogatz model \cite{watts1998collective} generates a ring lattice with $\lfloor m/n \rfloor$ neighbors per side, then rewires each edge with probability $\rho$. We evaluate three rewiring probabilities $\rho \in \{0.25, 0.50, 0.75\}$ and report the best.

\textbf{OURS.} A deterministic $O(n^2)$ algorithm that constructs a $k$-regular skeleton graph designed for high algebraic connectivity, described in Section 4.

\subsection{Results}

Figure~\ref{fig:comparison} presents the algebraic connectivity $\lambda_2$ achieved by each method as a function of the edge budget $m$, for graph sizes $n \in \{16, 32, 64, 128\}$. The REFINE agent is trained on small instances ($n \leq 16$) and evaluated without retraining at all sizes. At $n = 16$ (the largest training size), REFINE matches or exceeds the best baseline across the full density range. As $n$ increases beyond the training distribution, REFINE continues to outperform all baselines, demonstrating that the size-normalized features (Eq.~\ref{eq:features}) and the lightweight MLP policy generalize smoothly to graphs an order of magnitude larger than those seen during training.

\begin{figure}[t]
    \centering
    \includegraphics[width=\textwidth]{figure.png}
    \caption{Algebraic connectivity $\lambda_2$ vs.\ edge budget $m$ for $n \in \{16, 32, 64, 128\}$. Each subplot compares REFINE against the FV, ER, SW($\rho$), and OURS baselines. REFINE is trained on $n \leq 16$ and evaluated zero-shot at larger sizes.}\label{fig:comparison}
\end{figure}
\section{Ablation Studies}\label{sec:ablation}

A natural question is whether the learned policy is necessary, or whether the REFINE rewiring infrastructure alone—with a hand-crafted scoring rule—suffices. To isolate the contribution of RL, we compare two variants:

\begin{itemize}
    \item \textbf{REFINE (Greedy)}: The same iterative add/remove loop (Algorithm~\ref{alg:refine_inference}), but edges are scored using the raw Fiedler gap $(v_{2,i} - v_{2,j})^2$ instead of the learned MLP. This is equivalent to running the Fiedler vector heuristic \emph{within} the REFINE rewiring framework—the spectral subspace is still tracked via Rayleigh-Ritz, and bridge constraints are still enforced, but the policy is replaced by greedy first-order perturbation scoring.
    \item \textbf{REFINE (RL)}: The full method with the trained MLP policy ($\sim$13.6K parameters), as described in Section~\ref{sec:rr}.
\end{itemize}

Table~\ref{tab:ablation} reports the average $\lambda_2$ achieved by each variant across all $(n,m)$ configurations for each graph size. The greedy variant captures the benefit of the REFINE infrastructure (iterative rewiring, RR tracking, bridge safety), while the RL variant adds the benefit of \emph{learned} edge scoring.

\begin{table}[t]
\centering
\caption{Ablation: average $\lambda_2$ across all valid $(n, m)$ configurations. REFINE (Greedy) uses the Fiedler gap heuristic for edge scoring; REFINE (RL) uses the trained MLP policy. Higher is better.}\label{tab:ablation}
\begin{tabular}{@{}lcc@{}}
\toprule
$n$ & \textbf{REFINE (Greedy)} & \textbf{REFINE (RL)} \\
\midrule
16  & 6.131  & 6.663 \\
32  & 12.201  & 13.451 \\
64  & 24.840  &  27.694 \\
128 & 50.961  & 57.243  \\
\bottomrule
\end{tabular}
\end{table}

The results demonstrate that the REFINE infrastructure with greedy scoring provides only a modest improvement over the standalone FV baseline, as the hand-crafted heuristic makes the same myopic edge selections regardless of the rewiring context. The learned policy, by contrast, captures nonlinear interactions between spectral features and degree structure that the first-order perturbation approximation cannot express. This gap widens at larger $n$, precisely the regime where generalization matters most, confirming that the RL component is essential to REFINE's performance.


\section{Conclusion and Future Work}















\bibliographystyle{unsrt}  
\bibliography{references}




\end{document}
